
\SGASection{5}{Representability of the functor of smooth subgroups
equal to their connected normalizer}
\SGARefTarget{sga3:xv:section:5}

\SGAStatement{\SGAInlineRefTarget{sga3:xv:proposition:5-1}Proposition 5.1}{Let \(S\) be a prescheme, let \(G\)
be an \(S\)-group prescheme of finite presentation, and let \(H\) be
a subgroup prescheme of \(G\), smooth with connected fibres. Then:
\begin{enumerate}
\renewcommand{\labelenumi}{\textup{\alph{enumi})}}
\item the normalizer \(N\) of \(H\) in \(G\) is represented by a
closed subprescheme of \(G\), of finite presentation over \(S\);
\item\SGAInlineRefTarget{sga3:xv:proposition:5-1:item:b} the following conditions are equivalent:
\begin{enumerate}
\renewcommand{\labelenumii}{\textup{\roman{enumii})}}
\item the canonical immersion \(H\to N\) is an open immersion;
\item \(N\) is smooth along the identity section, and its identity
component, which is then representable (Exposé
VI\textsubscript{B}~\SGARefLink{sga3:xv:lemma:3:10}{3.10}), equals \(H\);
\item for every \(s\in S\), one has \(H_s=(N_s)^0\).
\end{enumerate}
\end{enumerate}}
\SGARefTarget{sga3:xv:proposition:5:1}

\SGAProof{Proof}
The group prescheme \(H\) is locally of finite presentation over
\(S\), since it is smooth, and has connected fibres; it is therefore
of finite presentation over \(S\) (Exposé
VI\textsubscript{B}~\SGARefLink{par:sga3-VIA-5.3.3}{5.3.3}). Part (a) follows from Exposé \SGARefLink{sga3:xi:corollary:6:11}{XI~6.11}. The
equivalence in \SGARefLink{sga3:xv:proposition:5-1:item:b}{(b)}, recalled here for reference, was proved in Exposé
VI\textsubscript{B}~\SGARefLink{prop:sga3-VIB-6.5.1}{6.5.1}.

We can now state the main theorem of this section.

\SGAStatement{\SGAInlineRefTarget{sga3:xv:theorem:5-2}Theorem 5.2}{Let \(S\) be a prescheme and \(G\) an
\(S\)-group prescheme of finite presentation over \(S\). Let
\(\mathcal L_G\), or simply \(\mathcal L\), be the \(S\)-functor
defined, for every \(S\)-prescheme \(S'\), by
\[
\mathcal L_G(S')=
\left\{
\begin{array}{c}
\text{subgroup preschemes of \(G_{S'}\), smooth over \(S'\),}\\
\text{with connected fibres, equal to their connected normalizer}
\end{array}
\right\}.
\]
Then \(\mathcal L\) is represented by an \(S\)-prescheme that is the
union of an increasing family \((U_i)_{i\in\mathbb N}\) of open
subpreschemes, each quasi-projective and of finite presentation over
\(S\), and hence in particular separated over \(S\).}
\SGARefTarget{sga3:xv:theorem:5:2}

\medskip\noindent\textbf{Initial reductions.}
For every \(r\geq0\), let \(\mathcal L_r\) be the subfunctor of
\(\mathcal L\) defined by
\[
\mathcal L_r(S')=
\left\{
\begin{array}{c}
H\subseteq G_{S'}:\ H\text{ is smooth with connected fibres,}\\
H\text{ equals its connected normalizer, and }\dim(H/S')=r
\end{array}
\right\}.
\]
The fibre dimension of a smooth group is locally constant. Hence
\(\mathcal L_r\to\mathcal L\) is represented by an immersion that is
both open and closed. It is therefore enough to prove the asserted
representability and finiteness properties for every
\(\mathcal L_r\), for then
\[
  \mathcal L=\coprod_{r\in\mathbb N}\mathcal L_r.
\]

For \(n\geq0\), an \(S\)-prescheme \(S'\), and a subgroup prescheme
\(H\) of \(G_{S'}\), denote by \(H^{(n)}\) the \(n\)-th normal
invariant of the identity section \(S'\to H\) (EGA IV~16.1.2).
Thus \(H^{(n)}\) is a finite-type \(\mathcal O_{S'}\)-module
corresponding to the \(n\)-th infinitesimal neighbourhood of the
identity section of \(H\). If \(H\) is smooth of relative dimension
\(r\), then \(H^{(n)}\) is locally free of a rank
\(\varphi(n,r)\) depending only on \(n,r\). Since \(H\) is a
subprescheme of \(G_{S'}\), there is a canonical epimorphism,
compatible with base change,
\[
 G^{(n)}\otimes_{\mathcal O_S}\mathcal O_{S'}
 \xrightarrow{\sim}G_{S'}^{(n)}
 \longrightarrow H^{(n)}.
\]

Introduce the projective \(S\)-scheme
\[
  P_{\varphi(n,r)}
   =\operatorname{Grass}_{\varphi(n,r)}(G^{(n)})
\]
(EGA I, 2nd ed., 9.7; see also Cartan Seminar 1960/61,
Exposé 14 by A. Grothendieck). The assignment
\[
  H\longmapsto H^{(n)}
\]
therefore defines a canonical morphism
\[
  u_{n,r}:\mathcal L_r\longrightarrow P_{\varphi(n,r)}.
\]
The group \(G\) acts naturally on \(G^{(n)}\), and hence on
\(P_{\varphi(n,r)}\), through
\[
 \operatorname{int}:G\longrightarrow
 \operatorname{Aut}_{S\text{-}\mathrm{gr}}(G),
 \qquad g\longmapsto\operatorname{int}(g).
\]

If \(S'\) is quasi-compact over \(S\) and \(H\in\mathcal L_r(S')\),
then for sufficiently large \(n\) one has (Exposé \SGARefLink{sga3:xi:corollary:6:11}{XI~6.11})
\[
 \operatorname{Norm}_{G_{S'}}(H)
 =\operatorname{Norm}_{G_{S'}}(H^{(n)}).
\]
For each \(n\geq0\), define the subfunctor \(\mathcal L_r^n\) of
\(\mathcal L_r\) by
\[
\mathcal L_r^n(S')=
\left\{
 H\in\mathcal L_r(S')\ \middle|\
 \operatorname{Norm}_{G_{S'}}(H)
 =\operatorname{Norm}_{G_{S'}}(H^{(n)})
\right\}.
\]

\medskip\noindent
\textbf{Representability of \(\mathcal L_r^n\).}
The integer \(r\) remains fixed through the proof, so we write
\(\mathcal L,\mathcal L^n,P_n,u_n\) in place of
\(\mathcal L_r,\mathcal L_r^n,P_{\varphi(n,r)},u_{n,r}\). Let
\[
 v_n:\mathcal L^n\longrightarrow P_n
\]
be the restriction of \(u_n\). The definition of \(\mathcal L^n\)
and \SGARefLink{sga3:xv:proposition:5-1:item:b}{Proposition 5.1(b)}(ii) show that \(v_n\) is a monomorphism.
More precisely:

\SGAStatement{Lemma 5.3}{The morphism \(v_n\) is an immersion of
finite presentation. Consequently \(\mathcal L^n\) is represented by
an \(S\)-prescheme, quasi-projective and of finite presentation over
\(S\).}
\SGARefTarget{sga3:xv:lemma:5:3}

After replacing \(S\) by \(P_n\), the usual method reduces the
assertion to the following. Let \(Q\in P_n(S)\), and define a
subfunctor \(F\) of the final functor \(h_S\) by
\[
F(S')=
\begin{cases}
h_S(S'),&\text{if there exists }H\in\mathcal L^n(S')
                    \text{ with }H^{(n)}=Q_{S'},\\
\varnothing,&\text{otherwise.}
\end{cases}
\]
We must prove that \(F\to S\) is an immersion of finite presentation.

Let \(N\) be the normalizer of \(Q\) in \(G\). It is a subgroup
prescheme of finite presentation over \(S\), namely the inverse image
of \(Q\) under
\[
  G\longrightarrow P_n,\qquad g\longmapsto g\cdot Q.
\]
We claim that \(F\) is equivalently described by
\[
F(S')=
\begin{cases}
h_S(S'),&
\begin{array}{l}
N_{S'}\text{ is smooth along the identity section, of}\\
\text{relative dimension }r,\text{ and }(N_{S'})^{(n)}=Q_{S'},
\end{array}\\[2mm]
\varnothing,&\text{otherwise.}
\end{cases}
\]

Call the first and second descriptions \(F_1\) and \(F_2\).
Suppose \(F_1(S')=h_S(S')\), and let \(H\in\mathcal L^n(S')\) with
\(H^{(n)}=Q_{S'}\). Then
\[
\begin{aligned}
\operatorname{Norm}_{G_{S'}}(H)
 &=\operatorname{Norm}_{G_{S'}}(H^{(n)})\\
 &=\operatorname{Norm}_{G_{S'}}(Q_{S'})
 =N_{S'}.
\end{aligned}
\]
By \SGARefLink{sga3:xv:proposition:5-1:item:b}{Proposition 5.1(b)}(ii), \(N_{S'}\) is smooth along the identity
section and its identity component is \(H\). Thus it has relative
dimension \(r\); since \(H\) is open in \(N_{S'}\),
\[
  (N_{S'})^{(n)}=H^{(n)}=Q_{S'}.
\]
Hence \(F_2(S')=h_S(S')\).

Conversely, suppose \(F_2(S')=h_S(S')\). The identity component of
\(N_{S'}\) is then represented (Exposé VI\textsubscript{B}~\SGARefLink{sga3:xv:lemma:3:10}{3.10}) by
a subgroup prescheme \(H\), smooth over \(S'\), with connected fibres
of dimension \(r\). Since \(H\) is invariant and open in \(N_{S'}\),
\[
\begin{aligned}
N_{S'}&\subseteq\operatorname{Norm}_{G_{S'}}(H)
\subseteq\operatorname{Norm}_{G_{S'}}(H^{(n)})\\
&=\operatorname{Norm}_{G_{S'}}((N_{S'})^{(n)})
=\operatorname{Norm}_{G_{S'}}(Q_{S'})=N_{S'}.
\end{aligned}
\]
All inclusions are equalities. The first says that \(H\) equals its
connected normalizer, and the second says
\(H\in\mathcal L^n(S')\). Hence \(F_1=F_2\).

We retain the second description. Condition (a) is represented by an
immersion of finite presentation, by the following lemma.

\SGAStatement{\SGAInlineRefTarget{sga3:xv:lemma:5-4}Lemma 5.4}{Let \(S\) be a prescheme, let \(X\) be
locally of finite presentation over \(S\), let \(\sigma:S\to X\) be
a section, let \(r\geq0\), and define the subfunctor \(L\) of \(S\)
by
\[
L(S')=
\begin{cases}
h_S(S'),&
\begin{array}{l}
X_{S'}\text{ is smooth along }\sigma_{S'},\text{ of relative}\\
\text{dimension }r\text{ at the points of }\sigma_{S'}(S'),
\end{array}\\[1mm]
\varnothing,&\text{otherwise.}
\end{cases}
\]
Then:
\begin{enumerate}
\renewcommand{\labelenumi}{\textup{\alph{enumi})}}
\item \(L\to S\) is an immersion of finite presentation;
\item let \(\mathcal J\) be the conormal sheaf of \(S\to X\).
If for every \(s\in S\) the vector space
\(\mathcal J\otimes_{\mathcal O_{S,s}}\kappa(s)\) has rank at most
\(r\), then \(L\to S\) is a closed immersion.
\end{enumerate}}
\SGARefTarget{sga3:xv:lemma:5:4}

\SGAProof{Proof}
The functor \(L\) is local on \(S\), so we may suppose \(S\) affine.
After replacing \(X\) by a neighbourhood of \(\sigma(S)\), we may
suppose \(X\) of finite presentation, and then reduce to noetherian
\(S\) by EGA IV~8.9. In part (b), use also compatibility of the
conormal sheaf with base change (EGA IV~16.6.4) and constructibility
of the rank of the fibres of \(\mathcal J\).

For an \(S\)-prescheme \(S'\), let
\(\mathcal J'=\mathcal J_{S'}\), let
\(S^\bullet(\mathcal J')\), canonically isomorphic to
\(S^\bullet(\mathcal J)_{S'}\), be the symmetric algebra of
\(\mathcal J'\) over \(\mathcal O_{S'}\), let
\(\operatorname{Gr}^\bullet(\sigma_{S'})\) be the graded algebra
associated with \(\sigma_{S'}\) (EGA IV~16), and let
\[
 \sigma_{n,S'}:S^n(\mathcal J')
 \longrightarrow\operatorname{Gr}^n(\sigma_{S'})
\]
be the canonical epimorphism. By EGA IV~17.12.3 and
EGA \(0_{\mathrm{IV}}\)~19.5.4, \(X_{S'}\) is smooth along
\(\sigma_{S'}\) of relative dimension \(r\) if and only if:
\begin{enumerate}
\renewcommand{\labelenumi}{\textup{\roman{enumi})}}
\item \(\mathcal J'\) is locally free of rank \(r\);
\item \(\sigma_{n,S'}\) is an isomorphism for every \(n\geq1\).
\end{enumerate}

TDTE IV, \SGARefLink{sga3:xv:lemma:3:6}{Lemma 3.6}, shows that the functor making \(\mathcal J\)
locally free of rank \(r\) is represented by a subprescheme \(S_1\),
closed in \(S\) in case (b). Replacing \(S\) by \(S_1\), assume that
\(\mathcal J\) is locally free of rank \(r\). Proceed by induction on
\(n\). Suppose a closed subprescheme \(S_{n-1}\) represents the
functor making \(\sigma_{q,S'}\) injective for \(q\leq n-1\).
The subfunctor\footnote{The editors replaced ``functor'' by
``subfunctor'' here, and ask whether the intended wording is instead
``which makes \(\sigma_{q,S'}\) injective for \(q\leq n\)''.}
making \(\sigma_{n,S'}\) injective is represented by a closed
subprescheme \(S_n\) of \(S_{n-1}\).

Indeed, after replacing \(S\) by \(S_{n-1}\), all
\(\sigma_{q,S}\) for \(q\leq n-1\) are isomorphisms. The
\((n-1)\)-st normal invariant \(X^{(n-1)}\) of \(\sigma_S\) has a
filtration with successive quotients
 \(\operatorname{Gr}^0(\sigma_S),\ldots,
 \operatorname{Gr}^{n-1}(\sigma_S)\), all locally free; hence
\(X^{(n-1)}\) is flat over \(S\). Compatibility of the normal
invariants with base change gives
\[
\begin{aligned}
\operatorname{Gr}^n(\sigma_{S'})
 &=\ker\!\left(X_{S'}^{(n)}\to X_{S'}^{(n-1)}\right)
  =\operatorname{Gr}^n(\sigma)_{S'},\\
\sigma_{n,S'}&=(\sigma_{n,S})_{S'}.
\end{aligned}
\]
Thus we must represent the functor making \(\sigma_{n,S}\) injective.
Locally we may suppose \(S\) affine and \(S^n(\mathcal J)\) free.
The required functor is represented by the closed subscheme defined
by the ideal generated by the coordinates of
\(\ker\sigma_{n,S}\) in a basis of \(S^n(\mathcal J)\).

Since \(S\) is noetherian, the descending chain \((S_n)\) stabilizes.
Its stable value represents \(L\), proving the lemma.

Return to \(\mathcal L^n\). After replacing \(S\) by the subprescheme
provided by \SGARefLink{sga3:xv:lemma:5-4}{Lemma 5.4}, assume \(N\) smooth along the identity section
of relative dimension \(r\). Then \(\mathcal L^n\) is the coincidence
functor of two sections \(h,g:S\to P_n\), corresponding to \(Q\) and
\(N^{(n)}\). It is therefore represented by the inverse image of the
diagonal of \(P_n\times_SP_n\) under the finite-presentation morphism
\(h\times_Sg\). This proves \SGARefLink{sga3:xv:lemma:5:3}{Lemma 5.3}.

\medskip\noindent
\textbf{The transition morphisms
\(\mathcal L^n\to\mathcal L^m\).}
If \(H\in\mathcal L^n(S')\), then \(H\in\mathcal L^m(S')\) for every
\(m\geq n\). Thus there are natural monomorphisms
\[
  u_n^m:\mathcal L^n\longrightarrow\mathcal L^m
  \qquad(m\geq n).
\]

\SGAStatement{Lemma 5.5}{The morphism \(u_n^m\) is an open immersion.}
\SGARefTarget{sga3:xv:lemma:5:5}

After a base change, it is enough to consider \(H\in\mathcal L^m(S)\)
and
\[
 N=\operatorname{Norm}_G(H)=\operatorname{Norm}_G(H^{(m)}),
 \qquad
 N'=\operatorname{Norm}_G(H^{(n)}).
\]
Let \(D\) be the coincidence functor of \(N,N'\): \(D(S')=h_S(S')\)
if \(N_{S'}=N'_{S'}\), and \(D(S')=\varnothing\) otherwise. We must
show that \(D\to S\) is an open immersion.

The same functor is obtained by requiring that the immersion
\(H\to N'\) become open. One implication follows from
\SGARefLink{sga3:xv:proposition:5-1}{Proposition 5.1}. Conversely, if \(H_{S'}\to N'_{S'}\) is open, then,
because \(H\) has connected fibres, \(H_{S'}\) is the identity
component of \(N'_{S'}\) (Exposé VI\textsubscript{B}~\SGARefLink{sga3:xv:lemma:3:10}{3.10}), hence is
invariant in \(N'_{S'}\). Thus \(N'_{S'}\subseteq N_{S'}\). The
opposite inclusion always holds, so equality follows. Since \(H,N'\)
are of finite presentation and \(H\) is flat over \(S\), Exposé
VI\textsubscript{B}~\SGARefLink{sga3:xv:remark:2:6}{2.6} shows that \(D\to S\) is open.

\medskip\noindent\textbf{End of the proof of \SGARefLink{sga3:xv:theorem:5-2}{Theorem 5.2}.}
The functors \(\mathcal L^n\) are represented, and the transition
morphisms \(u_n^m\) are compatible open immersions. There is therefore
an \(S\)-prescheme \(X\), the union of an increasing sequence of open
subpreschemes \((X_i)_{i\in\mathbb N}\), such that \(X_i\) represents
\(\mathcal L^i\) and \(X_i\to X_j\) corresponds to \(u_i^j\). In the
category of sheaves for the Zariski topology on \(\mathrm{Sch}/S\),
\[
  X=\varinjlim_iX_i,\qquad
  \mathcal L=\varinjlim_i\mathcal L^i.
\]
Thus \(X\) represents \(\mathcal L\), completing the proof.

\SGAStatement{Remark 5.6}{Suppose all residue characteristics of
\(S\) are zero. Then, for every \(r\geq0\),
\(\mathcal L_r=\mathcal L_r^1\), so \(\mathcal L_r\) is represented
by an \(S\)-prescheme quasi-projective and of finite presentation over
\(S\).}
\SGARefTarget{sga3:xv:remark:5:6}

Indeed, for \(H\in\mathcal L_r(S')\), it suffices to prove that
\[
  H\longrightarrow N=\operatorname{Norm}_{G_{S'}}(H^{(1)})
\]
is an open immersion. Because \(H\) is flat and both \(H,N\) are of
finite presentation, this may be checked on the fibres (Exposé
VI\textsubscript{B}~\SGARefLink{sga3:xv:remark:2:6}{2.6}). Cartier's theorem (Exposé
VI\textsubscript{B}~\SGARefLink{sga3:xv:proposition:1-6}{1.6}) gives, for a connected algebraic subgroup
\(H\) of an algebraic group in characteristic zero,
\[
 \operatorname{Norm}_G(H)
 =\operatorname{Norm}_G(\operatorname{Lie}H)
 =\operatorname{Norm}_G(H^{(1)}).
\]
\footnote{For example, one may apply Proposition II.6.1 of
M. Demazure and P. Gabriel, \emph{Groupes algébriques} I,
Masson--North Holland, 1970.}

In nonzero residue characteristic the subfunctors
\(\mathcal L_r^n\) can form a strictly increasing sequence, even when
\(S\) is quasi-compact, and then \(\mathcal L_r\) is not represented
by an \(S\)-prescheme quasi-compact over \(S\). For example, over a
field \(k\) of characteristic \(p>0\), let
\[
 T=\mathbb G_m\times\mathbb G_m,\qquad
 U=\mathbb G_a\times\mathbb G_a,
\]
and let \(G=T\ltimes U\), where
\((t,t')\cdot(u,u')=(tu,t'u')\). For \(n>0\), let \(U_n\subset U\)
be defined by \(u'=u^{p^n}\), and let \(T_n\subset T\) be defined by
\(t'=t^{p^n}\). Then \(T_n\) acts on \(U_n\), and
\(G_n=T_nU_n\) is smooth, connected, and equal to its normalizer in
\(G\). For every \(m\), the distinct groups \(G_n\), \(n\geq m\),
have the same infinitesimal neighbourhood of order \(p^m\).

\SGAStatement{Remark 5.7}{There is a canonical invertible sheaf
\(L\) on \(\mathcal L\) whose restriction to every open subset of
\(\mathcal L\) quasi-compact over \(S\) is \(S\)-ample. One may take
\[
  L=(\det(\operatorname{Lie}H))^{-1},
\]
where \(H\subseteq G_{\mathcal L}\) is the universal subgroup.}
\SGARefTarget{sga3:xv:remark:5:7}

Recall that for a finite-rank locally free sheaf \(\mathcal F\), its
determinant restricts on the open-and-closed locus where
\(\operatorname{rank}\mathcal F=r\) to \(\bigwedge^r\mathcal F\).
It is enough to restrict to \(\mathcal L_r^n\). Let
\(v_n:\mathcal L_r^n\to P_n\) be the canonical immersion, and let
\(Q\) be the universal locally free quotient on the Grassmannian
\(P_n\). By construction,
\[
  v_n^*Q=H^{(n)}.
\]
Since \(\det Q\) is the canonical ample sheaf on \(P_n\) (EGA I,
2nd ed., 9.7), \(\det H^{(n)}\) is ample relative to \(S\)
(EGA II~4.6.13(i) bis).

Let \(\mathcal J=H^{(1)}\) be the conormal sheaf, and let
\(S^q(\mathcal J)\) be the homogeneous degree-\(q\) part of its
symmetric algebra. Smoothness of \(H\) gives
\[
  \det H^{(n)}
  \xrightarrow{\sim}
  \bigotimes_{1\leq q\leq n}\det S^q(\mathcal J).
\]
For every locally free \(\mathcal J\) of rank \(r\) and every \(q>0\)
there is a canonical isomorphism
\[
  \det S^q(\mathcal J)
  \xrightarrow{\sim}
  (\det\mathcal J)^{\otimes s},
\]
where \(s>0\) depends only on \(r,q\). Thus, for some \(s>0\),
\[
  \det H^{(n)}\simeq(\det\mathcal J)^{\otimes s}.
\]
Since \(\det\mathcal J=(\det\operatorname{Lie}H)^{-1}\), EGA
II~4.5.6 proves that \(L\) is \(S\)-ample.

\SGAStatement{Remark 5.8}{Let \(G,H\) be \(S\)-group preschemes of
finite presentation and let \(i:H\to G\) be a monomorphism. Suppose
\(H\) is smooth over \(S\) with connected fibres. The normalizer
\(N=\operatorname{Norm}_G(H)\) is then represented by a closed
subgroup prescheme of \(G\), of finite presentation over \(S\)
(Exposé \SGARefLink{sga3:xi:corollary:6:11}{XI~6.11}). If, in addition, \(N\) is smooth along the identity
section and has the same relative dimension as \(H\), then
\(H=N^0\); in particular \(H\) is a subgroup prescheme of \(G\).}
\SGARefTarget{sga3:xv:remark:5:8}

Indeed, \(N^0\) is represented by an open subgroup prescheme of \(N\),
smooth over \(S\) (Exposé VI\textsubscript{B}~\SGARefLink{sga3:xv:lemma:3:10}{3.10}), and \(i\)
factors through \(N^0\). For every \(s\in S\), the connected smooth
algebraic groups \(H_s,(N^0)_s\) have the same dimension and hence
coincide. Since \(H\) is flat, \(H\to N^0\) is an isomorphism
(EGA IV~17.9.5). Thus the functor \(\mathcal L\) of this section may
equivalently be described as the functor of subgroup preschemes
\(H\) of \(G\) that are smooth over \(S\), have connected fibres,
and equal their connected normalizer.
